% ============================================================
%  Informe 2 Práctica 2 — SOII
%  Compilar: lualatex informe2_practica2.tex  (dos pasadas)
% ============================================================
\documentclass[a4paper, 12pt]{article}

% ── Codificación y fuentes ──────────────────────────────────
\usepackage[T1]{fontenc}
\usepackage[utf8]{inputenc}

% ── Idioma ──────────────────────────────────────────────────
\usepackage[spanish, es-nolists, es-noindentfirst]{babel}

% ── Geometría ───────────────────────────────────────────────
\usepackage[
  a4paper,
  top=2.5cm, bottom=2.5cm,
  left=2.8cm, right=2.8cm
]{geometry}

% ── Colores ─────────────────────────────────────────────────
\usepackage{xcolor}
\definecolor{azul}     {HTML}{2E5FA3}
\definecolor{azulosc}  {HTML}{1A3A6B}
\definecolor{grisclaro}{HTML}{F4F4F4}
\definecolor{grisborde}{HTML}{CCCCCC}
\definecolor{rojocod}  {HTML}{A31515}

% ── Encabezados y pies ──────────────────────────────────────
\usepackage{fancyhdr}
\pagestyle{fancy}
\fancyhf{}
\fancyhead[L]{\small\color{azul}\textbf{SOII --- Práctica 2}}
\fancyhead[R]{\small\color{gray}Informe 2: Sincronización con semáforos}
\fancyfoot[C]{\small\thepage}
\renewcommand{\headrulewidth}{0.4pt}

% ── Títulos de sección ──────────────────────────────────────
\usepackage{titlesec}
\titleformat{\section}
  {\Large\bfseries\color{azul}}
  {\thesection.}{0.6em}{}
  [\vspace{-0.3em}\textcolor{azul}{\rule{\linewidth}{0.8pt}}\vspace{0.2em}]

\titleformat{\subsection}
  {\large\bfseries\color{azulosc}}
  {\thesubsection.}{0.5em}{}

\titleformat{\subsubsection}
  {\normalsize\bfseries}
  {\thesubsubsection.}{0.4em}{}

% ── Listas ──────────────────────────────────────────────────
\usepackage{enumitem}
\setlist[itemize]{leftmargin=1.5em, itemsep=2pt, topsep=4pt}

% ── Tablas ──────────────────────────────────────────────────
\usepackage{booktabs}
\usepackage{array}
\usepackage{colortbl}
\newcolumntype{L}[1]{>{\raggedright\arraybackslash}p{#1}}
\newcolumntype{C}[1]{>{\centering\arraybackslash}p{#1}}

% ── Código fuente ───────────────────────────────────────────
\usepackage{listings}
\lstdefinestyle{cstyle}{
  language=C,
  basicstyle=\ttfamily\small,
  keywordstyle=\bfseries\color{azul},
  commentstyle=\itshape\color{gray},
  stringstyle=\color{rojocod},
  numberstyle=\tiny\color{gray},
  numbers=left,
  numbersep=8pt,
  stepnumber=1,
  backgroundcolor=\color{grisclaro},
  frame=single,
  framesep=4pt,
  rulecolor=\color{grisborde},
  breaklines=true,
  tabsize=4,
  showstringspaces=false,
  xleftmargin=1.8em,
  xrightmargin=0.5em,
  captionpos=b,
  literate=
    {á}{{\'a}}1 {é}{{\'e}}1 {í}{{\'i}}1 {ó}{{\'o}}1 {ú}{{\'u}}1
    {Á}{{\'A}}1 {É}{{\'E}}1 {Í}{{\'I}}1 {Ó}{{\'O}}1 {Ú}{{\'U}}1
    {ñ}{{\~n}}1 {Ñ}{{\~N}}1,
}
\lstdefinestyle{shellstyle}{
  basicstyle=\ttfamily\small,
  backgroundcolor=\color{grisclaro},
  frame=single,
  framesep=4pt,
  rulecolor=\color{grisborde},
  breaklines=true,
  numbers=none,
  xleftmargin=1em,
  xrightmargin=0.5em,
}
\lstset{style=cstyle}

% ── Cajas resaltadas ────────────────────────────────────────
\usepackage{tcolorbox}
\tcbuselibrary{skins, breakable}

\newtcolorbox{infobox}[1][]{
  colback=azul!7, colframe=azul, fonttitle=\bfseries,
  title=#1, left=6pt, right=6pt, top=4pt, bottom=4pt,
  breakable
}
\newtcolorbox{warnbox}[1][]{
  colback=orange!8, colframe=orange!70!black, fonttitle=\bfseries,
  title=#1, left=6pt, right=6pt, top=4pt, bottom=4pt,
  breakable
}

% ── Hiperenlaces ────────────────────────────────────────────
\usepackage{hyperref}
\hypersetup{
  colorlinks=true, linkcolor=azul, urlcolor=azul,
  pdftitle={Práctica 2 SOII --- Informe 2},
}

% ── Miscelánea ──────────────────────────────────────────────
\usepackage{microtype}
\usepackage{parskip}
\setlength{\parskip}{6pt}

% ── Macro inline-code ───────────────────────────────────────
\newcommand{\ic}[1]{\texttt{#1}}

% ============================================================
\begin{document}

% ── Portada ─────────────────────────────────────────────────
\begin{titlepage}
  \centering
  \vspace*{3cm}
  {\Huge\bfseries\color{azul} Sistemas Operativos II\par}
  \vspace{0.6cm}
  {\LARGE Práctica 2 --- Informe 2\par}
  \vspace{0.4cm}
  {\large\color{gray} Sincronización de procesos con semáforos\par}
  \vspace{0.5cm}
  \textcolor{azul}{\rule{0.7\linewidth}{1.2pt}}
  \vspace{2cm}

  \begin{tabular}{rl}
    \textbf{Asignatura:} & Sistemas Operativos II \\[4pt]
    \textbf{Titulación:} & Grado en Ingeniería Informática \\[4pt]
    \textbf{Curso:}      & 2024/25 \\[4pt]
    \textbf{Fecha:}      & \today \\
  \end{tabular}

  \vfill
  {\small Universidad de Santiago de Compostela}
\end{titlepage}

\tableofcontents
\newpage

% Para que la sección comience numerada como "2"
\setcounter{section}{1}

% ============================================================
\section{Informe 2 --- Solución con semáforos POSIX}
% ============================================================

\subsection{Descripción de la solución}

La solución clásica al problema del productor-consumidor utiliza tres
semáforos POSIX que abordan las tres causas de error del apartado anterior:

\vspace{0.5em}
\begin{center}
\begin{tabular}{L{2.4cm} C{2.2cm} L{8.5cm}}
  \toprule
  \rowcolor{azul!80}
  \color{white}\textbf{Semáforo} &
  \color{white}\textbf{Val.\ inicial} &
  \color{white}\textbf{Propósito} \\
  \midrule
  \ic{VACIAS}  & \ic{N = 10} &
    Número de huecos libres. El productor hace \ic{sem\_wait} antes
    de insertar; se bloquea si el buffer está lleno. \\
  \addlinespace
  \ic{LLENAS}  & \ic{0} &
    Número de elementos disponibles. El consumidor hace \ic{sem\_wait}
    antes de retirar; se bloquea si el buffer está vacío. \\
  \addlinespace
  \ic{MUTEX}   & \ic{1} &
    Exclusión mutua en la región crítica (acceso a \ic{buf} y
    \ic{top}). \\
  \bottomrule
\end{tabular}
\end{center}

\subsection{Cómo ejecutar los programas}

\begin{lstlisting}[style=shellstyle]
gcc -o prod_sem prod_sem.c -pthread
gcc -o cons_sem cons_sem.c -pthread

# Terminal 1:
./prod_sem
# Terminal 2:
./cons_sem
\end{lstlisting}

El productor crea e inicializa los semáforos con \ic{O\_CREAT}.
El consumidor los abre con \ic{sem\_open("NOMBRE", 0)},
reintentando hasta que el productor los haya creado.

\subsection{Protocolo de sincronización}

\subsubsection*{Productor}

\begin{lstlisting}[style=cstyle,
  caption={Esquema del bucle del productor (\texttt{prod\_sem.c})}]
for (int iter = 0; iter < MAX_ITER; iter++) {
    char item = produce_item();   // FUERA de la region critica

    sem_wait(vacias);   // bloquear si buffer lleno
    sem_wait(mutex);    // entrar en region critica
    /* ====== REGION CRITICA ====== */
    insert_item(item);
    int val_top = shm->top; // capturar top dentro de la RC
    /* ============================ */
    sem_post(mutex);    // salir de region critica
    sem_post(llenas);   // avisar al consumidor

    sleep(fase_actual); // FUERA de la region critica
}
\end{lstlisting}

\subsubsection*{Consumidor}

\begin{lstlisting}[style=cstyle,
  caption={Esquema del bucle del consumidor (\texttt{cons\_sem.c})}]
for (int iter = 0; iter < MAX_ITER; iter++) {
    sem_wait(llenas);   // bloquear si buffer vacio
    sem_wait(mutex);    // entrar en region critica
    /* ====== REGION CRITICA ====== */
    char item = remove_item();
    int val_top = shm->top; // capturar top dentro de la RC
    /* ============================ */
    sem_post(mutex);    // salir de region critica
    sem_post(vacias);   // avisar al productor

    consume_item(item); // FUERA de la region critica
    sleep(fase_actual); // FUERA de la region critica
}
\end{lstlisting}

\begin{infobox}[Orden crítico de las llamadas]
  \ic{sem\_wait(condición)} debe preceder siempre a
  \ic{sem\_wait(mutex)}.
  Si se invirtiese el orden, un proceso podría retener el \ic{MUTEX}
  mientras espera en el semáforo de condición, bloqueando al otro
  proceso indefinidamente (\emph{interbloqueo}).
\end{infobox}

\subsection{Fases de velocidad}

El bucle de 80 iteraciones tiene tres fases con \ic{sleep()} situados
\textbf{fuera} de la región crítica:

\vspace{0.5em}
\begin{center}
\begin{tabular}{C{2.2cm} C{2.2cm} C{2.2cm} L{5.5cm}}
  \toprule
  \rowcolor{azul!80}
  \color{white}\textbf{Iteraciones} &
  \color{white}\textbf{Prod.} &
  \color{white}\textbf{Cons.} &
  \color{white}\textbf{Efecto esperado} \\
  \midrule
  0 -- 29 & 0\,s & 2\,s &
    Buffer lleno. Productor bloqueado en \ic{sem\_wait(vacias)}
    cuando \ic{top = N}. \\
  \addlinespace
  30 -- 59 & 2\,s & 0\,s &
    Buffer vacío. Consumidor bloqueado en \ic{sem\_wait(llenas)}
    cuando \ic{top = 0}. \\
  \addlinespace
  60 -- 79 & rand 0--3\,s & rand 0--3\,s &
    Comportamiento mixto; \ic{top} oscila entre 0 y~4. \\
  \bottomrule
\end{tabular}
\end{center}

\subsection{Resultados observados}

La salida de \ic{prod\_sem} muestra \ic{top\_tras\_insert = 10}
de forma sostenida durante las iteraciones 10--39, lo que confirma que el
productor avanza sin esperar al consumidor (fase~1) y el buffer se mantiene
lleno. En la fase~2 (iter 30--59) \ic{top\_tras\_insert} cae a~1,
evidenciando que el consumidor vacía el buffer casi en tiempo real.
En la fase aleatoria los valores de \ic{top} oscilan entre~1 y~4,
reflejando el comportamiento no determinista de los tiempos aleatorios.

El recuento final de vocales fue idéntico en productor y consumidor:

\vspace{0.5em}
\begin{center}
\begin{tabular}{C{3cm} C{3cm} C{3cm}}
  \toprule
  \rowcolor{azul!80}
  \color{white}\textbf{Vocal} &
  \color{white}\textbf{Producidas} &
  \color{white}\textbf{Consumidas} \\
  \midrule
  a\,/\,A & 8 & 8 \\
  e\,/\,E & 6 & 6 \\
  i\,/\,I & 8 & 8 \\
  o\,/\,O & 5 & 5 \\
  u\,/\,U & 3 & 3 \\
  \bottomrule
\end{tabular}
\end{center}

\subsection{Gestión de semáforos en el kernel}

Los semáforos POSIX nombrados persisten en el kernel hasta que se eliminan
explícitamente. Para evitar conflictos con ejecuciones anteriores, el
productor llama a \ic{sem\_unlink()} \emph{antes} de crear los semáforos y
de nuevo \emph{al terminar}:

\begin{lstlisting}[style=cstyle]
/* Al inicio: eliminar posibles residuos de ejecuciones previas */
sem_unlink("VACIAS"); sem_unlink("LLENAS"); sem_unlink("MUTEX");

/* Crear e inicializar */
vacias = sem_open("VACIAS", O_CREAT, 0700, N);
llenas = sem_open("LLENAS", O_CREAT, 0700, 0);
mutex  = sem_open("MUTEX",  O_CREAT, 0700, 1);

/* ... bucle productor ... */

/* Al terminar: cerrar y eliminar del kernel */
sem_close(vacias); sem_close(llenas); sem_close(mutex);
sem_unlink("VACIAS"); sem_unlink("LLENAS"); sem_unlink("MUTEX");
\end{lstlisting}

El consumidor únicamente cierra sus descriptores con \ic{sem\_close()},
sin eliminar los semáforos del kernel.

\subsection{Conclusiones}

\begin{itemize}
  \item Los semáforos POSIX resuelven de forma correcta y eficiente las
        tres causas de error del apartado anterior: \ic{MUTEX} garantiza
        exclusión mutua; \ic{VACIAS} bloquea al productor sin
        \emph{busy-waiting}; \ic{LLENAS} bloquea al consumidor sin
        \emph{busy-waiting}.
  \item El orden de llamadas \ic{sem\_wait(condición)} $\to$
        \ic{sem\_wait(mutex)} es fundamental para evitar interbloqueos.
        Invertirlo provocaría que un proceso retuviese el mutex mientras
        espera en el semáforo de condición, bloqueando al otro proceso
        indefinidamente.
  \item Las fases de velocidad demuestran visualmente que los semáforos
        de condición actúan como señales \emph{sleep}/\emph{wakeup}: el
        proceso lento queda suspendido sin consumir CPU hasta que el otro
        proceso libera un hueco o deposita un elemento.
  \item La coincidencia exacta de los conteos de vocales en todas las
        ejecuciones es la prueba objetiva de que no se ha perdido ni
        duplicado ningún carácter, independientemente de los
        \emph{timings} aleatorios de la fase~3.
  \item La función \ic{sem\_unlink()} es imprescindible para una
        correcta limpieza de recursos del kernel; su omisión deja
        semáforos residuales que interfieren con ejecuciones posteriores.
\end{itemize}

\end{document}